\section{Introduction}\label{sec:intro}

The subword tokenizer determines the string unit on which a language model operates.
Common tokenizers like byte-pair encoding~\cite{sennrich2016bpe} and the unigram language model~\cite{kudo2018subword} learn the boundaries of unit tokens from the statistical co-occurrence of characters alone.
It is for this reason that they can never learn that \emph{walking} is composed of \emph{walk} and \emph{ing}, that \emph{the} usually precedes nouns and adjectives, or that Turkish \emph{evlerimizden} happily deposits a cascade of four suffixes at the end of \emph{ev}.
In recent years there has been a growing effort to incorporate additional morphological knowledge into tokenizers~\cite{jabbar2023morphpiece,asgari2025morphbpe,hofmann2021superbizarre}, and in this work we explore the capabilities of adding syntactic knowledge.
POS taggers for more than 60 languages are widely available~\cite{qi2020stanza}, and tagging a corpus is a one-shot process.
POS tags provide a reusable signal about the grammar of a text and are unified across languages under the Universal Dependencies scheme~\cite{nivre2020ud}.

The three tokenizers we evaluate use the POS tags from the tagger as a signal and their lemmas to guide training on a BPE vocabulary.
The first method we describe weighs the BPE merge probability based on the likelihood of a POS transition across word boundaries.
The second weights the probability of merges inside content words so that nouns, verbs, and adjectives stay whole.
The third, called Lemma2Word, uses the lemma from a POS tagger to produce a boundary indicating the stem of a word which it then segments before subword training.
We also evaluate a large language model prompted to segment words, as a segmenter rather than a tokenizer.

Initial tests without held-out words, gold boundaries, or repeated samples showed promising results on the core tokenizer evaluation metrics (segmentation entropy, token length, and vocabulary stability) but did not directly address the use of a tokenizer.
In the controlled evaluation that follows, we provide that evidence for 5 languages: English, Spanish, Danish, Tamil, and Turkish.
Methods were trained with the same Wikipedia text under an \nVocabBudget-type budget and each was repeated three times using different samples to train the tokenizers.
We compared all methods to both BPE and the unigram model across a series of metrics: boundary precision and recall against gold boundaries produced by MorphScore~\cite{arnett2025morphscore} (derived from treebank lemmas), compression and stability on held-out text, and bits-per-character for a small language model trained on the output of each tokenizer.

The first method that weights BPE merges based on POS transition probability hardly affects the results, since it acts only on tokens that absorb a space at their ends.
It produces the same pairs of frequent function words as unweighted BPE and changes the boundary F1 less than the variance between seed-to-seed runs for four of the five languages.
The second method that weights merges inside content words affects each language differently.
The improvement in Tamil is substantial but it is actually a negative effect for Spanish and Turkish.
The third method, Lemma2Word, increases boundary recall by 23--34 points for four languages and by 77 points in Tamil in a way that adds 4--8\% more tokens.
Across English, Tamil, and Turkish, Lemma2Word retains most of this gain even on words its lexicon has never seen.

Finally, we conclude that the unigram language model remains the best-performing tokenizer that makes no use of linguistic knowledge and is the only one competitive in Tamil without a lemmatizer.
The language model probe does not follow the boundary scores and in all five languages ranks the unigram model as first and Lemma2Word as last in bits per character.
In other words, finding more morpheme boundaries does not guarantee a better tokenizer for a language model, at least at the 7M scale we tested.
All code, data preparation scripts, and numbers are produced by the released repository.
